\chapter{No algorithm can distinguish the types}\label{tx:sec:undecidability}
Already in the special case of the fixed target $G_n=1$, the
classification cannot be turned into an algorithm
which determines the case from a sentence. This remains true under the
promise that every finite instance is inconsistent. Throughout this section,
the input is a well-formed normalised sentence and the request is a
refutation of its translated axioms, with the translation of
Chapter~\ref{tx:sec:systems}; the characteristic $p$ is fixed.
For a sentence $\psi$, write $d_S(\psi,n)$ for the least degree of
deriving $1$ from its size-$n$ translation. This is its refutation
profile. Its type is the alternative in Theorem~\ref{tx:thm:classification}:
eventually bounded (T1), no recurring finite level (T2), not eventually bounded with
a finite nonempty set of recurring levels (T3), or infinitely many
recurring finite levels (T4).

\section{Knowing the alternatives need not identify the alternative}
The classification gives four exhaustive structural patterns. It does
not give an algorithm deciding which pattern belongs to a presented
uniform family. The reduction below makes this difference sharp even
for the constant target one, where every finite instance is inconsistent.

Think of a program whose halting is unknown. A uniform computation-prefix
sentence has models of every finite size if the program never halts.
If it does halt, a fixed finite prefix supplies a bounded-degree
contradiction at all large sizes. We use this switch inside a controlled
join of two example families. The join has one type in the nonhalting
case and the other type in the halting case. An algorithm distinguishing
those types would therefore decide halting.

Three algebraic points need proof, and each is supplied below. The
halting prefix yields an actual bounded-degree NS certificate, not just
a semantic contradiction. Adjoining a satisfiable block preserves the
other branch's degree. The guarded disjunction keeps both branches
visible with at most an additive degree-one overhead. These are the
connections between the computational reduction and our exact proof
systems.

An algorithm receives a finite string encoding a sentence: its symbols,
clauses and function arities can all be read effectively. \emph{Partial
computable} allows the algorithm to run forever on some inputs.
The theorem demands termination and a correct answer only on the
\emph{promised} inputs, namely sentences with no finite nonempty model
and type in $A\cup B$. It does not demand a procedure for recognising
that promise. The reduction will produce promised inputs in both
the halting and the nonhalting case, so a separator could not evade
them by failing to halt.

\begin{theorem}\label{tx:thm:undecidability}
Let $A,B$ be disjoint nonempty subsets of $\{\mathrm T_1,\mathrm T_2,
\mathrm T_3,\mathrm T_4\}$. There is no partial computable algorithm which,
on every normalised sentence having no finite nonempty model and having
type in $A\cup B$, halts and correctly distinguishes membership in $A$
from membership in $B$. This holds for NS and for ordinary PC, over every
field of characteristic $p$.
\end{theorem}

In particular, each of the six pairs is undecidable, and so is membership
in any nonempty proper union of types. The last conclusion is the scope
of the historical every-union theorem. The pairwise statement requires the
additional construction below; it does not follow merely by listing all
the undecidable union problems.

\section{A sentence that records a computation prefix}
A \emph{reduction} transforms each finite program description into a
finite sentence description by a fixed effective procedure. It must
not first decide whether the program halts. Calling the transformation
\emph{primitive recursive} states a stronger effectiveness property:
it is built using composition and recursion whose iteration counts
are given as natural-number arguments. The later unbounded search is performed by the
program, not by the procedure producing its sentence.

A counter program has finitely many natural-number registers and instructions
for incrementing a register, branching on zero and decrementing, and halting.
All registers initially contain zero. The number of registers is finite
but need not be two. Write $c_M$ for one more than the largest register index
occurring in $M$, so $c_M\ge1$.

From a program $M$, form a sentence $\theta_M$ with a strict linear order,
least and greatest elements $o,e$, a successor function $S$, counter-value
functions $C_r$, and state predicates $Q_\ell$. Require $S(e)=e$ and require
$S(t)$ to be the immediate successor of $t$ when $t\ne e$. At $o$, all
counters equal $o$ and the state is the initial one. Exactly one state holds
at each element, and halting states are forbidden.

At $t\ne e$, the state determines the transition to $S(t)$. Increment
replaces $C_r(t)$ by $S(C_r(t))$, with the guard $C_r(t)\ne e$.
Decrement branches according as $C_r(t)=o$; on the positive branch, the new
value $v$ satisfies $v<C_r(t)$ and $S(v)=C_r(t)$. Other counters are
unchanged. These are finite universal clauses with flattened function
graphs, and $M\mapsto\theta_M$ is primitive recursive.

\begin{lemma}[Computation switch]\label{tx:lem:switch}
If $M$ never halts, $\theta_M$ has a model of every size $n\ge1$.
If the run of $M$ first reaches a halting state at time $T$, the translated system has an NS refutation of
degree bounded by a constant $C_M$ at every size $n\ge T+1$.
\end{lemma}
\begin{proof}
In the first case, use $[n]=\{0,\ldots,n-1\}$ in its natural order and
record the first $n$ configurations. Every counter value at time $j$ is at
most $j$, so no guarded increment overflows before the last element.
This includes $n=1$.

For the second case put $a_j=S^j(o)$. If the structure has $T+1$ distinct
elements, then $a_j\ne e$ for $j<T$: otherwise immediacy of successor and
linearity of the order leave only $a_0,\ldots,a_j$. Induction along this
finite prefix forces the actual state at $a_j$ and forces
$C_r(a_j)=a_{c_r(j)}$, where $c_r(j)$ is the numerical counter value.
At $a_T$ this gives a forbidden halting state. The argument needs neither
well-ordering nor an induction axiom in the structure.

This local contradiction gives a degree bound by the finite-term
certificate in Lemma~\ref{tx:lem:finite-certificate}. More explicitly, with $c_M$
counters one may take
\[
 m=T+2+2c_M(T+1),\qquad
 C_M=m+\sum_R(m+T+1)^{\operatorname{arity}(R)}
                  +\max_{C\text{ a clause}}|C|.
\]
The term nodes are $o,e,a_1,\ldots,a_T,C_r(a_j),S(C_r(a_j))$;
$T+1$ distinct additional labels supply the size condition. The constant
depends on the program and its halting time, never on $n$.
\end{proof}

For completeness, the reduction needs effective universality, not a
decision procedure for the halting time. Every primitive recursive function
can be compiled into finite counter code: composition uses fresh registers,
and primitive recursion uses a counted loop, preserving the input registers.
Let $H(k,c)$ equal one when the program with code $c$, on its fixed
encoded input, halts within $k$ simulated steps, and zero otherwise.
Bounded simulation computes $H$ by a fixed number of primitive-recursive
transition updates, so this predicate is primitive recursive. A counter program that successively tests
$H(0,c),H(1,c),\ldots$ therefore halts exactly when the coded program does.
Loading $c$ into this fixed search program is primitive recursive. This
reduces the usual undecidable halting problem to the programs just used.
Only effective compilation and the standard undecidability of halting are used; the number of counters is unrestricted but finite.

\section{Finite local contradictions and their certificates}
\label{tx:app:finite-certificate}
The halting switch uses a fixed finite computation prefix even though
function totality is expressed by sums with $n$ summands. The following
elementary certificate construction separates these two sizes.

An acyclic evaluation list is an ordered list of term evaluations
whose arguments have already been named. This is a dependency condition
on the list, not a prohibition on a function taking the same value twice;
constants can be treated as evaluations with no arguments. The proof
sums over every labelling of that finite list. A \emph{falsity indicator}
for a Boolean literal is $1-x$ for a positive literal and $x$ for a
negative literal. The translated clause is the product of its literal
falsity indicators: it vanishes exactly when at least one literal is
true. These products are the actual input polynomials used below.

\begin{lemma}[Finite-term certificate]\label{tx:lem:finite-certificate}
Let $n\ge1$ and let a finite acyclic list of $m$ function evaluations be given, with every
argument supplied by an earlier node. Fix $b$ distinct labels. Suppose that
every assignment of labels to the nodes, and every Boolean assignment to
the relation atoms on the named and fixed labels, either has inconsistent
function values or falsifies a clause instance whose variables lie among the named and fixed
labels and all of whose function evaluations occur in the list. If clauses have length at most $L$, the
graph-and-sum translation has an NS refutation of degree at most
\[
                 m+\sum_R(m+b)^{\operatorname{arity}(R)}+L.
\]
The hypothesis and conclusion concern sizes for which the $b$ distinct
labels exist.
\end{lemma}
\begin{proof}
For a labelling $\alpha$ of the nodes, let $q_\alpha$ be the product of the
$m$ graph variables asserting those function evaluations. Expanding the
totality sum for each node, in the order of the acyclic list, gives
\[
                         \sum_\alpha q_\alpha-1
\]
in the degree-$m$ NS span. This is a telescoping identity: at step $j$, the
product for the first $j-1$ evaluations multiplies the totality equation
for evaluation $j$. Summing over the possible earlier values gives the
next partial sum. Repeated graph variables cause no problem.

If $\alpha$ gives different values to the same function on the same input,
$q_\alpha$ contains a functionality collision and already belongs to the
span within degree $m$. Otherwise, list the distinct relation variables
on the at most $m+b$ named labels. Write this actual set as
$Y_\alpha=\{y_1,\ldots,y_{a_\alpha}\}$, where
$a_\alpha\leq a_{\mathrm{rel}}=\sum_R(m+b)^{\operatorname{arity}(R)}$.
Expand over these distinct variables:
\[
             1=\prod_{i=1}^{a_\alpha}(y_i+(1-y_i))
              =\sum_\beta u_\beta.
\]
Each $u_\beta$ is the Boolean assignment monomial for one partial relational
structure. By hypothesis that structure falsifies an available clause.
Its clause polynomial is a product of at most $L$ falsity indicators.
Write $P=q_\alpha u_\beta$ and the clause polynomial as
$F_1\cdots F_\ell$, where each $F_j$ is the falsity indicator of a
literal false in the chosen partial structure. For each literal,
$P(1-F_j)$ has an explicit expression in the structural NS span:
\begin{itemize}
\item For a positive relation literal, $F_j=1-x$ and $u_\beta=(1-x)u'$.
Thus $P(1-F_j)=q_\alpha u'x(1-x)=-q_\alpha u'(x^2-x)$.
For a negative relation literal, $F_j=x$ and $u_\beta=xu'$; the same
identity with $x(1-x)$ applies.
\item For a negative graph literal false in the partial structure,
its graph variable $x$ is one of the chosen evaluations, so
$q_\alpha=xq'$. Its falsity indicator is $F_j=x$ and
$P(1-F_j)=-q'u_\beta(x^2-x)$.
\item For a positive graph literal false there, the chosen evaluation
has a different output variable $x'$ on the same function input.
Here $q_\alpha=x'q'$, $F_j=1-x$, and
$P(1-F_j)=q'u_\beta x'x$, an explicit functionality-axiom multiple.
\item A false equality or disequality has the constant falsity
indicator $F_j=1$, so $P(1-F_j)=0$.
\end{itemize}
These are all cases of the normalised literals. Each identity has
product degree at most $m+a_{\mathrm{rel}}+1$.
The exact telescoping identity
\[
 P(1-F_1\cdots F_\ell)
   =\sum_{j=1}^{\ell}P(1-F_j)\prod_{i<j}F_i
\]
therefore puts its left side in the NS span within degree
$m+a_{\mathrm{rel}}+\ell$. Since $PF_1\cdots F_\ell$ is a
multiple of the falsified clause axiom, $P$ is in that same span.
Sum first over $\beta$, then over $\alpha$, and use the totality identity
to obtain one.
\end{proof}


A \emph{partial structure} for $\theta_M$ has a domain of labels, all relations defined on it, and
function values defined only on some arguments. It is \emph{locally consistent} if every clause
instance whose variables take values in the domain, and whose function atoms are all defined,
has a true literal.

\begin{lemma}[Local switch]\label{tx:lem:local-switch}
Let the run of $M$ first reach a halting state at time $T$, and let $M$ have $c_M$ counters.
Suppose a locally consistent partial structure defines $o$, $e$, $a_{j+1}=S(a_j)$ for $j<T$
\textup{(}with $a_0=o$\textup{)}, $C_r(a_j)$ and $S(C_r(a_j))$ for $j\le T$, $r<c_M$, and that all these
elements lie in its domain. Then its domain has at most $T$ distinct elements.
\end{lemma}
\begin{proof}
Suppose it contains $T+1$ distinct elements. If $a_j=e$ for some $j<T$, take $j$ least. Every
domain element $x$ satisfies $x=a_0$ or $a_0<x$ (least element), and, if $a_i<x$ with $i<j$, then
$x=a_{i+1}$ or $a_{i+1}<x$ (no element strictly between $a_i$ and $S(a_i)$, and linearity). As
$a_j=e$ is greatest, the domain is $\{a_0,\ldots,a_j\}$, too small. So $a_j\ne e$ for $j<T$, and
the order clauses make $a_0<a_1<\cdots<a_T$. For the counter induction, the initial clauses give the initial
state and $C_r(a_0)=a_0$ for every counter. A counter value at time $j$
is at most $j$, so the notation $a_v$ uses only this named prefix.
Suppose its states and counters agree through time $j<T$.
An increment of value $v$ gives $S(a_v)=a_{v+1}$; because $v\le j<T$,
this successor is among the defined evaluations and its guard is valid.
The zero branch of a decrement is selected exactly when $v=0$.
On a positive decrement, the transition specifies a value $y<a_v$
with $S(y)=a_v$. The immediate-predecessor property forces
$y=a_{v-1}$: if $y<a_{v-1}$ then $a_{v-1}$ lies strictly between
$y$ and $S(y)$, and if $a_{v-1}<y<a_v$ then $y$ lies strictly
between $a_{v-1}$ and $S(a_{v-1})$. Both alternatives contradict
immediacy. Unchanged counters retain their named values. State
uniqueness and the chosen transition identify the next state.
Every clause used has its variables in the domain and its function
atoms in the specified finite evaluation list. This proves the
induction through time $T$. At $j=T$ the state is halting, which a clause forbids.
\end{proof}

\begin{corollary}\label{tx:cor:switch-degree}
In that case, for every $n\ge T+1$ the translation of $\theta_M$ at size $n$ has an NS refutation
of degree at most
\[
 C_M=m+\sum_R(m+T+1)^{\operatorname{arity}(R)}+\max_{C}|C|,\qquad m=T+2+2c_M(T+1).
\]
\end{corollary}
\begin{proof}
Apply Lemma~\ref{tx:lem:finite-certificate} to the $m$ evaluations $o$, $e$, $a_1,\ldots,a_T$,
$C_r(a_j)$, $S(C_r(a_j))$ ($j\le T$, $r<c_M$), each argument supplied by an earlier node, with
$b=T+1$ fixed labels. A consistent labelling and a Boolean assignment give a partial structure of
the kind in Lemma~\ref{tx:lem:local-switch}, whose domain contains the $T+1$ fixed labels; so it is not
locally consistent, which is the hypothesis of Lemma~\ref{tx:lem:finite-certificate}.
\end{proof}

Here $c_M$ is one more than the largest counter index occurring in $M$ (so $c_M\ge1$), and $T$ must be
the first halting time. Sizes $n\le T$ are genuine exceptions: $\theta_M$ can have models there.

\section{A disjunction with controlled degree}
The notation $d_*(\alpha,n)$ below always refers to the original
graph-and-sum translation of $\alpha$, with target $1$; the symbol $*$
stands for either NS or ordinary PC.
A signature is the collection of relation and function symbols.
Disjoint signatures are obtained by renaming the symbols of one branch,
so the two polynomial variable families are independent even though
they use the same universe. \emph{Structural equations} below mean the
Boolean, function-totality and function-functionality equations of the
compiler. The clause equations impose the sentence's further demands.
This distinction lets the inactive branch retain a structural root
while its clause equations are guarded away.

Conjunction is easy to control when one block is satisfiable. In disjoint
signatures, if $\beta$ has a root at size $n$, then
\begin{equation}\label{tx:eq:satisfiable-block}
 d_*(\alpha\wedge\beta,n)=d_*(\alpha,n),
 \qquad *=\NS,\PC.
\end{equation}
One inequality follows by using the refutation of $\alpha$; the other
follows by substituting the root of $\beta$ into every proof line.

We also need a disjunction which does not accidentally simplify either
branch. Introduce a fresh unary relation $Z$. Keep both branches' Boolean,
totality and functionality equations. Replace every clause polynomial $f$
of $\alpha$ by $z_t f$, and every clause polynomial $g$ of $\beta$ by
$(1-z_t)g$, for a fresh universally quantified label $t$. Denote this
normalised sentence by $\alpha\join\beta$.
Thus setting every $z_t$ to one exposes the $\alpha$ clauses, while
setting them to zero exposes the $\beta$ clauses. Keeping the structural
equations makes either specialisation an actual copy of the selected
branch's polynomial system.

The guard $z_t$ is itself a Boolean relation variable. Multiplying a
clause by it makes the clause active when $z_t=1$; multiplying by
$1-z_t$ selects the other value. The lower-bound specialisations set
\emph{all} guards to the same value. For the upper bound one label
$t_0$ suffices, since the universally instantiated inputs include the
clauses guarded by that label. The hypothesis $n\ge1$ guarantees that
such a label exists.

\begin{lemma}[Controlled join]\label{tx:lem:join}
For $n\ge1$, with $\infty$ allowed, each of NS and PC satisfies
\[
 \max\{d_*(\alpha,n),d_*(\beta,n)\}
 \le d_*(\alpha\join\beta,n)
 \le\max\{d_*(\alpha,n),d_*(\beta,n)\}+1.
\]
\end{lemma}
\begin{proof}
Set all selectors to one, and give the inactive block a structural root:
all its functions take one fixed value and all its relations are false.
Its clauses are guarded away, while its structural equations vanish.
The resulting proof refutes $\alpha$ without increasing degree. Selectors
zero give the other lower inequality.

For the upper bound fix a label $t_0$. Simulate each branch refutation
with its fixed guard $c=z_{t_0}$ or $c=1-z_{t_0}$, representing a branch
line $f$ by $cf$. A guarded clause is already an axiom. A structural
input $h$ is retained unguarded, so PC can form $z_{t_0}h$ and,
for the other guard, $h-z_{t_0}h$. Thereafter a linear combination of
branch lines is simulated by the same combination of their guarded
lines, and a step $f\mapsto x_vf$ is simulated by the legal variable
multiplication $cf\mapsto x_v(cf)$. This uses no unguarded intermediate
branch line and remains valid after cancellations. Every line has
degree at most one above the branch budget. The two simulations end
in $z_{t_0}$ and $1-z_{t_0}$, whose sum is one. For NS the same
construction multiplies the two certificates by their guards.
\end{proof}

\section{The reduction}
Let $\psi_1,\psi_2,\psi_3,\psi_4$ be the four sentences of
Section~\ref{tx:sec:examples}. Each has no finite nonempty model. For these particular
witnesses, the dominant type under the maximum operation is ordered by
\[
                 \mathrm T_1\prec\mathrm T_4\prec\mathrm T_3\prec\mathrm T_2.
\]
This is a property of the chosen witnesses, not an ordering assertion about
arbitrary profiles. A bounded profile has no effect on the type of a maximum;
the all-size linear lower bound for $\psi_2$ makes it dominant; and the
maximum of $\psi_4$ and $\psi_3$ is bounded off the multiples of $p$ and
linear on them. An additive error of at most one preserves these conclusions.
More generally, adjoining an eventually bounded profile and allowing this
error preserves each of the four intrinsic types: recurrence transfers
through finite intervals of degrees, and unboundedness transfers directly.

Choose indices $a,b$ with $\mathrm T_a\prec\mathrm T_b$ in this order
and define, using disjoint signatures,
\[
             \Phi_M=\psi_a\join(\psi_b\wedge\theta_M).
\]
If $M$ does not halt, \eqref{tx:eq:satisfiable-block} makes the second branch
have exactly the profile of $\psi_b$. The join has type $\mathrm T_b$. If $M$ halts,
the second branch is eventually bounded by Lemma~\ref{tx:lem:switch}, and the
join has type $\mathrm T_a$. Every $\Phi_M$ is within the all-finite-inconsistent
promise, whether or not $M$ halts. The map is primitive recursive and
uses no information about that alternative. A separator of types $\mathrm T_a,\mathrm T_b$
would decide halting. Choosing $a$ and $b$ on opposite sides of any
disjoint nonempty $A,B$ proves Theorem~\ref{tx:thm:undecidability}.

\section{Where the textbook argument ends}
The general consequence theorem has reduced the target question to a
support-bounded closure calculation, fulfilled its transfer contract by
one common projector, and read off the four profiles from fixed-degree
recurrence. The refutation examples and this reduction show that the
four alternatives are substantial and cannot be algorithmically
separated in general. The exercise solutions which follow return to
small calculations and help check that each logical distinction in this
route has been retained.
